% Overleaf-ready ICCD experiment text.
% Required preamble package if not already present:
% \usepackage{booktabs}
%
% Recommended paste location:
% sections/5_experiments.tex, after "Computation-Communication Overlap Analysis".

\subsection{Supplemental Validation}
\label{sec:supplemental_validation}

\textbf{Goal and protocol.}
We add a controlled supplemental evaluation to answer four reviewer-facing
questions: whether verified prefetching preserves target outputs, whether the
gain survives fair cache-controlled baselines, whether online self-correction
introduces measurable overhead, and where SPICE saturates under denser routing.
Unless otherwise stated, the controlled harness uses 16 MoE layers, 64 routed
experts, Top-$K{=}6$, a 512-entry expert cache, and 8\,MB expert transfers. All
methods use the same cache capacity and host-to-device transfer model.

\textbf{Verified correctness.}
SPICE speculates expert residency, not target computation. In a verified MoE
execution harness, SPICE and on-demand offloading produce identical logits:
the maximum logit difference is 0.0, the exact argmax match rate is 100\%, and
the baseline and SPICE verified pseudo-perplexities are both 1035.21. The
prefetch-slot hit rate is 74.24\%, while the remaining 25.76\% slots are served
by fallback loading. This confirms that fallback affects latency rather than
model semantics.

\begin{table}[t]
\centering
\caption{Controlled offloading and prefetch comparison. All policies use the
same cache capacity and expert-transfer model. Lower TPOT and fallback are
better; higher hit rate indicates fewer critical-path fetches.}
\label{tab:supp_prefetch}
\scriptsize
\begin{tabular}{lrrrr}
\toprule
Policy & TPOT (ms) & Hit (\%) & Fallback (\%) & H2D (GB) \\
\midrule
Naive & 71.25 & 0.00 & 100.00 & 384.00 \\
LRU & 44.67 & 85.04 & 14.96 & 57.44 \\
MoE-Offloading & 44.67 & 85.04 & 14.96 & 57.44 \\
Pre-gated & 41.83 & 94.15 & 5.85 & 147.27 \\
\textbf{SPICE} & \textbf{41.03} & \textbf{99.76} & \textbf{0.24} & 198.80 \\
\bottomrule
\end{tabular}
\end{table}

\textbf{Main comparison.}
Table~\ref{tab:supp_prefetch} shows that SPICE reduces critical-path fallback
traffic from 5.85\% for Pre-gated to 0.24\%. The total H2D volume is higher
because SPICE issues more aggressive speculative transfers, but these transfers
are scheduled to overlap with target execution rather than becoming synchronous
fallback stalls.

\begin{table}[t]
\centering
\caption{Draft-path overhead and LoRE ablation. Online self-correction is
reported separately from the default offline draft path.}
\label{tab:supp_overhead}
\scriptsize
\begin{tabular}{lrrr}
\toprule
Variant & TPOT (ms) & Fallback (\%) & Online overhead \\
\midrule
\textbf{SPICE offline} & \textbf{41.04} & \textbf{0.24} & 0.00 ms \\
SPICE online & 42.96 & 0.24 & 983.04 ms \\
SPICE without LoRE & 42.56 & 5.10 & 0.00 ms \\
\bottomrule
\end{tabular}
\end{table}

\textbf{Overhead and ablation.}
Table~\ref{tab:supp_overhead} shows that online self-correction is measurable:
enabling online updates increases TPOT from 41.04\,ms to 42.96\,ms in this
setting. We therefore use the offline draft path as the default latency
configuration and treat online adaptation as an optional robustness mode. The
no-LoRE variant increases fallback from 0.24\% to 5.10\%, indicating that the
low-rank state correction is important for accurate verified prefetching.

\begin{table*}[t]
\centering
\caption{High-utilization Top-$K$ stress study with all baselines. The setting
uses 768 decoding steps, 64 experts, 8\,MB expert transfers, and the same tight
cache budget for every policy. Lower TPOT and fallback are better.}
\label{tab:supp_topk_baselines}
\scriptsize
\begin{tabular}{rrrrrrr}
\toprule
Top-$K$ & Naive & LRU & MoE-Offload & Pre-gated & SPICE & SPICE fallback (\%) \\
\midrule
2 & 50.42 & 42.92 & 42.92 & 41.66 & \textbf{41.01} & \textbf{0.46} \\
4 & 60.83 & 45.91 & 45.91 & 43.30 & \textbf{42.51} & \textbf{7.43} \\
6 & 71.25 & 49.10 & 49.10 & 45.09 & \textbf{43.14} & \textbf{6.97} \\
8 & 81.67 & 52.65 & 52.65 & 54.37 & \textbf{51.11} & \textbf{24.36} \\
10 & 92.08 & \textbf{56.25} & \textbf{56.25} & 68.90 & 58.36 & 33.41 \\
12 & 102.50 & \textbf{58.39} & \textbf{58.39} & 76.54 & 65.35 & 39.03 \\
\bottomrule
\end{tabular}
\end{table*}

\textbf{Scalability boundary.}
Table~\ref{tab:supp_topk_baselines} directly addresses the high-utilization
case: the Naive baseline already drives modeled PCIe activity from 20.66\% at
$K{=}2$ to 60.98\% at $K{=}12$, and Pre-gated reaches 34.91--100\% activity for
$K{\ge}4$. SPICE remains fastest through $K{=}8$ by reducing synchronous
fallback traffic, but the advantage disappears once dense routing pushes the
transfer demand beyond the overlap window. At $K{=}10$ and $K{=}12$, the
conservative LRU/MoE-Offloading cache has lower TPOT than SPICE because further
speculation consumes bandwidth that cannot be hidden. This supports a bounded
systems claim: SPICE improves offloaded MoE serving when useful overlap windows
exist, but it does not remove the physical host-to-device bandwidth limit.

\begin{table}[t]
\centering
\caption{Controlled proxy comparison with recent MoE inference systems. Proxy
rows are same-harness approximations, not official reproductions.}
\label{tab:supp_proxy}
\scriptsize
\begin{tabular}{lrrr}
\toprule
Variant & TPOT (ms) & Hit (\%) & Fallback (\%) \\
\midrule
ExpertFlow proxy & 41.70 & 98.66 & 1.34 \\
AdapMoE proxy & 42.08 & 94.38 & 5.62 \\
SP-MoE proxy & 43.04 & 95.39 & 4.61 \\
MoE-SpeQ proxy & 41.98 & 95.72 & 4.28 \\
\textbf{SPICE verified} & \textbf{41.03} & \textbf{99.78} & \textbf{0.22} \\
\bottomrule
\end{tabular}
\end{table}

\textbf{Comparison to recent designs.}
Table~\ref{tab:supp_proxy} isolates the design effect of routing-only verified
prefetching against controlled proxy variants inspired by recent MoE inference
systems. These results should not be read as official reproductions. Within the
same cache and PCIe model, SPICE has the lowest fallback rate because its
speculative decision controls expert residency and is verified against target
routing before expert execution.

\textbf{Power and timeline telemetry.}
We also collect GPU telemetry to avoid inferring energy behavior from latency
alone. In a copy-and-compute workload, measured GPU power averages 184.32\,W
and peaks at 452.98\,W over a 44.90\,s sampling window. The workload transfers
16\,GB in 0.569\,s, corresponding to an effective 28.13\,GB/s. We also collect
an Nsight Systems CUDA trace for the same workload. These measurements provide
runtime telemetry for copy/compute overlap, but we do not claim full
system-level energy savings without paired baseline-vs-SPICE energy runs.
